\documentclass[12pt]{article}
\input{iati_structure.tex}
\renewcommand{\bottomtitlespace}{3cm}

\usepackage[english]{babel}

\begin{document}

\renewcommand{\headerLang}{Română}
\renewcommand{\problemName}{AB13. Vision}

\renewcommand{\headerLeft}{IATI Day 1, Senior group}
\renewcommand{\headerRightFirst}{XVII INTERNATIONAL ADVANCED TOURNAMENT IN INFORMATICS}
\renewcommand{\headerRightSecond}{ONLINE 2026}

\renewcommand{\tl}{$5$ s}
\renewcommand{\ml}{$1024$ MB}
\problem{PROBLEMĂ \problemName}
\addauthor{Author: Emil Indzhev}{2026}{04}{16}

%The USS Enterprise spaceship (with Sashka as its new captain) is traveling throughout the universe, which is an infinite 2D grid of sectors. In each sector $(X, Y)$\footnote{$X$ is row number (increasing in the down direction) and $Y$ is the column number (increasing in the right direction).} there is a beacon with vision strength $V_{X, Y}$. When in that sector, the ship uses the beacon to scan all sectors $(X', Y')$ with $X - V_{X, Y} \leq X' \leq X + V_{X, Y}$ and  $Y - V_{X, Y} \leq Y' \leq Y + V_{X, Y}$. Thus, the seen sectors form a square of side length $2V_{X, Y} + 1$. For each such sector, the only information seen is the vision strength of the beacon in it, i.e. $V_{X', Y'}$ (this is because beacons act as both transmitters and receivers). After scanning, the ship then teleports to one of these seen sectors, as beacons also act as teleporters.

Nava spațială USS Enterprise (cu Sashka ca nou căpitan) călătorește prin univers, care este o grilă 2D infinită de sectoare. În fiecare sector $(X, Y)$\footnote{$X$ este numărul rândului (crescător în direcția în jos) și $Y$ este numărul coloanei (crescător în direcția dreaptă).} există un far cu o putere vizuală $V_{X, Y}$. Când se află în acel sector, nava folosește farul pentru a scana toate sectoarele $(X', Y')$ cu $X - V_{X, Y} \leq X' \leq X + V_{X, Y}$ și $Y - V_{X, Y} \leq Y' \leq Y + V_{X, Y}$. Astfel, sectoarele vizibile formează un pătrat cu latura de lungime $2V_{X, Y} + 1$. Pentru fiecare astfel de sector, singura informație vizibilă este puterea vizuală a farului din acesta, adică $V_{X', Y'}$ (acest lucru se datorează faptului că farurile acționează atât ca emițătoare, cât și ca receptoare). După scanare, nava se teleportează către unul dintre aceste sectoare vizibile, deoarece farurile acționează și ca teleportoare.

%Unfortunately, the beacons have one critical flaw: they may flip the $X$ axis of the scan, or the $Y$ axis, or both (i.e. there are $4$ possible scan orientations). Notice that the ship uses this possibly flipped scan to decide which sector to teleport to, but the teleportation also uses the same orientation, e.g. if the $Y$ axis is flipped and the ship teleports to the left direction (decreasing $Y$) of the scan, it will actually teleport to the right. This means that the ship will indeed teleport to sector chosen in the scan.

Din păcate, farurile au un defect critic: pot inversa axa $X$ a scanării, sau axa $Y$, sau ambele (adică există $4$ orientări posibile de scanare). Observați că nava folosește această scanare posibil inversată pentru a decide în ce sector să se teleporteze, dar teleportarea folosește și ea aceeași orientare, de exemplu, dacă axa $Y$ este inversată și nava se teleportează în direcția stângă a scanării ($Y$ mai mic), astfel se va teleporta de fapt în dreapta. Aceasta înseamnă că nava se va teleporta într-adevăr în sectorul ales în scanare.

%Furthermore, the ship has no memory, so its decision on where to teleport to must depend only on the scan produced by the beacon. Additionally, the ship's strategy must be deterministic -- if given the same scan, it must choose to teleport to the same sector in the scan.

În plus, nava nu are memorie, așa că decizia sa privind locul în care se va teleporta trebuie să depindă doar de scanarea produsă de semnalul farului. În plus, strategia navei trebuie să fie deterministă - dacă i se dă aceeași scanare, trebuie să aleagă să se teleporteze în același sector din scanare.

%The ship starts at unknown coordinates and must always (regardless of where it starts and regardless of what orientations of scans it gets) be able to travel in the direction of increasing $X$ and $Y$ arbitrarily far (we can say that we require it to eventually reach a sector $(X, Y)$ such that $X, Y \geq 10^{100}$). However, for this to be possible, both the travel strategy of the ship and the layout of vision strengths across the universe are crucial.

Nava pornește de la coordonate necunoscute și trebuie întotdeauna (indiferent de unde pornește și indiferent de orientările scanărilor pe care le primește) să fie capabilă să se deplaseze în direcția creșterii $X$ și $Y$ la o distanță arbitrară (putem spune că avem nevoie ca aceasta să ajungă în cele din urmă la un sector $(X, Y)$ astfel încât $X, Y \geq 10^{100}$). Totuși, pentru ca acest lucru să fie posibil, atât strategia de deplasare a navei, cât și configurația puterilor vizuale în univers sunt cruciale.

%This is where you come in -- you have to design an $N$ by $M$ rectangular pattern $P$ of vision strengths which will repeat infinitely across the universe: $V_{X, Y} = P_{X \bmod N, Y \bmod M}$. You also have to design the travel strategy for the ship, i.e. a function which takes in the scan produced by the beacon and chooses to which of those sectors to teleport to next.

Aici intervii tu - trebuie să proiectezi un model dreptunghiular $P$ de dimensiuni $N$ pe $M$, de puteri vizuale,  care se va repeta la infinit în univers: $V_{X, Y} = P_{X \bmod N, Y \bmod M}$. De asemenea, trebuie să proiectezi strategia de deplasare a navei, adică o funcție care primește scanarea produsă de far și alege în care dintre aceste sectoare să se teleporteze în continuare.

%Since beacons with higher vision strength are expensive, your goal is to produce a valid solution, while trying to minimize the average vision strength in your pattern, i.e. the average value of $P$.

Deoarece farurile cu o putere vizuală mai mare sunt scumpe, obiectivul este să produceți o soluție validă, încercând în același timp să minimizați puterea vizuală medie din modelul dvs., adică valoarea medie a lui $P$.

%Since this problem may seem quite challenging, you decide to start with a practice problem: the same problem but in 1D. In this simpler problem, we throw away the $X$ dimension and only use the $Y$ one. Thus, your pattern is just a sequence of length $M$ and it is repeated like so: $V_Y = P_{Y \bmod M}$. The scans produced by beacons are also just sequences of length $2V_Y + 1$ (containing the sectors $Y'$ such that  $Y - V_{X, Y} \leq Y' \leq Y + V_{X, Y}$). Here, there are only 2 possible orientations for the scan (regular or flipped) and the ship must be able to reach $Y \geq 10^{100}$.

Întrucât această problemă poate părea destul de dificilă, decideți să începeți cu o problemă mai simplă: aceeași problemă, dar în 1D. În această problemă mai simplă, eliminăm dimensiunea $X$ și o folosim doar pe cea $Y$. Astfel, modelul dvs. este doar o secvență de lungime $M$ și se repetă astfel: $V_Y = P_{Y \bmod M}$. Scanările produse de faruri sunt, de asemenea, doar secvențe de lungime $2V_Y + 1$ (conținând sectoarele $Y'$ astfel încât $Y - V_{X, Y} \leq Y' \leq Y + V_{X, Y}$). Aici, există doar 2 orientări posibile pentru scanare (cea normală și cea inversată) și nava trebuie să poată ajunge la $Y \geq 10^{100}$.

%Your solutions for the 1D case and the 2D case will be scored separately and you may get points for one without having a valid solution for the other.

Soluțiile dvs. pentru cazul 1D și cazul 2D vor fi punctate separat și este posibil să obțineți puncte pentru unul fără a avea o soluție validă pentru celălalt.

%\subsection{Implementation details}
\subsection{Detalii de implementare}

%You have to implement 4 functions: \verb|getVisionPattern1d|, \verb|getMove1d|, \verb|getVisionPattern2d| and \verb|getMove2d|. The first 2 are used for the 1D case and the second 2 -- for the 2D one. In any test only one of these pairs of functions will be called, but all 4 must be implemented (even if the implementations for one pair are empty) for the solution to be valid (to compile).

Trebuie să implementați 4 funcții: \verb|getVisionPattern1d|, \verb|getMove1d|, \verb|getVisionPattern2d| și \verb|getMove2d|. Primele 2 sunt folosite pentru cazul 1D, iar celelalte 2 -- pentru cel 2D. În orice test, va fi apelată doar una dintre aceste perechi de funcții, dar toate cele 4 trebuie implementate (chiar dacă implementările pentru o pereche sunt goale) pentru ca soluția să fie validă (să se compileze).

\begin{lstlisting}
 std::vector<int> getVisionPattern1d()
\end{lstlisting}

%This function will be called once, if the test is for the 1D case. It must return a the sequence $P$ of length $M$ -- the 1D vision strength pattern to be repeated. $M$ and all elements of $P$ must be between $1$ and $60$.

Această funcție va fi apelată o singură dată, dacă testul este pentru cazul 1D. Trebuie să returneze secvența $P$ de lungime $M$ -- modelul 1D al puterilor vizuale care trebuie repetat. $M$ și toate elementele lui $P$ trebuie să fie între $1$ și $60$.

\begin{lstlisting}
 int getMove1d(std::vector<int> v)
\end{lstlisting}

%This function will be called once per possible scan, if the test is for the 1D case. It is given a scan produced by a beacon, as a sequence of length $2V + 1$ (where $V$ is the vision strength of the beacon in the central sector in the sequence). It must return which sector the ship should teleport to, as an index $T$ in the sequence, such that $0 \leq T < 2V + 1$.

Această funcție va fi apelată o singură dată pentru o scanare posibilă, dacă testul este pentru cazul 1D. I se dă o scanare produsă de un far, ca o secvență de lungime $2V + 1$ (unde $V$ este puterea vizuală a farului din sectorul central din secvență). Trebuie să returneze sectorul în care nava ar trebui să se teleporteze, ca un index $T$ din secvență, astfel încât $0 \leq T < 2V + 1$.

\begin{lstlisting}
 std::vector<std::vector<int>> getVisionPattern2d()
\end{lstlisting}

%This function will be called once, if the test is for the 2D case. It must return a rectangular table $P$ of size $N$ by $M$ -- the 2D vision strength pattern to be repeated. $N$, $M$ and all elements of $P$ must be between $1$ and $60$.

Această funcție va fi apelată o singură dată, dacă testul este pentru cazul 2D. Trebuie să returneze un tabel dreptunghiular $P$ de dimensiune $N$ pe $M$ -- modelul de putere vizuală 2D care trebuie repetat. $N$, $M$ și toate elementele lui $P$ trebuie să fie între $1$ și $60$.

\begin{lstlisting}
 std::pair<int, int> getMove2d(std::vector<std::vector<int>> v)
\end{lstlisting}

%This function will be called once per possible scan, if the test is for the 2D case. It is given a scan produced by a beacon, as a square table of size $2V + 1$ by $2V + 1$ (where $V$ is the vision strength of the beacon in the central sector in the square). It must return which sector the ship should teleport to, as a pair of indices $S$ and $T$ in the table, such that $0 \leq S, T < 2V + 1$.

Această funcție va fi apelată o dată pentru fiecare scanare posibilă, dacă testul este pentru cazul 2D. I se oferă o scanare produsă de un far, sub forma unui tabel pătrat de dimensiunea $2V + 1$ pe $2V + 1$ (unde $V$ este puterea vizuală a farului din sectorul central din pătrat). Trebuie să returneze sectorul în care nava ar trebui să se teleporteze, ca o pereche de indici $S$ și $T$ din tabel, astfel încât $0 \leq S, T < 2V + 1$.

%For the given dimensionality $D$ ($1$ or $2$) of the test, the grader will check that your program produces a valid pattern -- that it makes valid teleportation choices for all possible scans and that your solutions works regardless of starting coordinates and sequence of orientations of scans.

Pentru dimensionalitatea dată $D$ ($1$ sau $2$) a testului, evaluatorul va verifica dacă programul dvs. produce un model valid - dacă face alegeri valide de teleportare pentru toate scanările posibile și dacă soluțiile dvs. funcționează indiferent de coordonatele de pornire și secvența orientărilor scanărilor.

\subsection{Restricții}

\vspace{0.1em}

\begin{itemize}
\item $1 \leq D \leq 2$
\item $1 \leq N, M, V_{X, Y}, V_Y \leq 60$
\end{itemize}

\newpage

\subsection{Subtask-uri și punctare}

\begin{table}[H]
\begin{tblr}{|Q[c,m]|Q[c,m]|Q[c,m]|Q[c,m]|}
\hline
\textbf{Subtask} & \textbf{Puncte} & \textbf{$D$} & \textbf{$A^*$} \\
\hline
$1$ & $24$ & $1$ & $1.5$ \\
\hline
$2$ & $76$ & $2$ & $1.125$ \\
\hline
\end{tblr}
\end{table}
\FloatBarrier

%Your score for a given subtask (if your solution for it is correct) depends on the average vision strength in your pattern $P$ for it. Let that be $A$ and let the target for the subtask be $A^*$. Your score (fraction of points for the subtask) will be:

Scorul tău pentru un anumit subtask (dacă soluția ta este corectă) depinde de puterea vizuală medie în modelul tău $P$ pentru subtask. Fie aceasta $A$ și ținta pentru subtask $A^*$. Scorul tău (fracțiunea de puncte pentru subtask) va fi:

$$1 - {\left(1 - \frac{A^* - 1}{\max{\left(A, A^*\right)} - 1}\right)}^{0.8}$$

%\subsection{Sample pattern}
\subsection{Exemplu de model}

%Suppose your solution returns the following pattern with $N=3$ and $M=2$:

Să presupunem că soluția ta returnează următorul model cu $N=3$ și $M=2$:

\[
\left[
\begin{array}{c c}
1 & 2 \\
3 & 4 \\
5 & 6
\end{array}
\right]
\]

%Now, let's examine the possible scans at sector $(0, 1)$, which has vision strength $2$. There are $4$ possible orientations (some of which produce the same scan on this pattern):

Acum, să examinăm posibilele scanări la sectorul $(0, 1)$, care are o putere vizuală de $2$. Există $4$ orientări posibile (unele dintre ele produc aceeași scanare pe acest model):

\[
\begin{array}{cc}
%\text{Orientation $0$: No flips} & \text{Orientation $1$: $Y$-axis flipped} \\
\text{Orientare $0$: Fără inversări} & \text{Orientare $1$: Axa $Y$ inversată} \\
\left[
\begin{array}{ccccc}
4 & 3 & 4 & 3 & 4 \\
6 & 5 & 6 & 5 & 6 \\
2 & 1 & 2 & 1 & 2 \\
4 & 3 & 4 & 3 & 4 \\
6 & 5 & 6 & 5 & 6
\end{array}
\right]
&
\left[
\begin{array}{ccccc}
4 & 3 & 4 & 3 & 4 \\
6 & 5 & 6 & 5 & 6 \\
2 & 1 & 2 & 1 & 2 \\
4 & 3 & 4 & 3 & 4 \\
6 & 5 & 6 & 5 & 6
\end{array}
\right]
\\
\\
%\text{Orientation $2$: $X$-axis flipped} & \text{Orientation $3$: Both axes flipped} \\
\text{Orientare $2$: axa $X$ inversată} & \text{Orientare $3$: Ambele axe inversate} \\
\left[
\begin{array}{ccccc}
6 & 5 & 6 & 5 & 6 \\
4 & 3 & 4 & 3 & 4 \\
2 & 1 & 2 & 1 & 2 \\
6 & 5 & 6 & 5 & 6 \\
4 & 3 & 4 & 3 & 4
\end{array}
\right]
&
\left[
\begin{array}{ccccc}
6 & 5 & 6 & 5 & 6 \\
4 & 3 & 4 & 3 & 4 \\
2 & 1 & 2 & 1 & 2 \\
6 & 5 & 6 & 5 & 6 \\
4 & 3 & 4 & 3 & 4
\end{array}
\right]
\end{array}
\] 

%Since $0$ and $1$ are identical and since the ship is deterministic, only one call to \verb|getMove2d| will be made for them. Suppose your solution returns the move $(0, 1)$ on this scan. On orientation $0$ this would mean moving $1$ to the left and $2$ up, while on orientation $1$ -- moving $1$ to the right and $2$ up.


Deoarece orientările $0$ și $1$ sunt identice și deoarece nava este deterministă, se va face un singur apel către \verb|getMove2d| pentru aceste orientări. Să presupunem că soluția ta returnează teleportarea $(0, 1)$ în această scanare. La orientarea $0$, aceasta ar însemna mutarea cu $1$ la stânga și cu $2$ în sus, în timp ce la orientarea $1$ -- mutarea cu $1$ la dreapta și cu $2$ în sus.

%\subsection{Sample grader}
\subsection{Structura Evaluării}

%The sample grader's first input is $D$. After that it will do all calls to your function: first getting the pattern $P$ and then caching all teleportation choices. After that, it will start a console interaction, asking for the starting coordinates of the ship. After that, it will print the current state: coordinates and ``raw'' (without flips) scan. Then, it will ask for the orientation to be used ($0$ for no flips, $1$ to flip the $Y$ axis, $2$ to flip the $X$ axis and $3$ to flip both). After that, it will print the scan in the given orientation, as well as what move the ship choose. The ship will be moved and this will repeat. Notice that the sample grader will not automatically verify that your solution is correct.

Primul input al evaluatorului este $D$. După aceea, va efectua toate apelurile către funcția dvs.: mai întâi va obține modelul $P$ și apoi va memora în cache toate opțiunile de teleportare. După aceea, va începe o interacțiune cu consola, solicitând coordonatele inițiale ale navei. După aceea, va afișa starea curentă: coordonate și scanarea „brută” (fără inversări). Apoi, va solicita orientarea care va fi utilizată ($0$ pentru fără inversări, $1$ pentru a inversa axa $Y$, $2$ pentru a inversa axa $X$ și $3$ pentru a inversa ambele). După aceea, va afișa scanarea în orientarea dată, precum și mișcarea aleasă de navă. Nava va fi mutată și acest lucru se va repeta. Observați că evaluatorul nu va verifica automat dacă soluția dvs. este corectă.

\end{document}
